% bank_fintech_tone/drafts/zh/manuscript.tex
% P4.2 中文 manuscript 骨架 (plan-r3 衍生论文扩展, 2026-05-15)
% 编译: xelatex + biber + xelatex × 2

% papers_assets/_shared/preamble_zh.tex
% Shared Chinese LaTeX preamble (XeLaTeX + CJKutf8).
% P4.1 (plan-r3 derivative-paper extension, 2026-05-15)
% Compile: xelatex + biber + xelatex × 2

\documentclass[12pt,a4paper]{article}

% --- CJK Fonts (XeLaTeX, cross-platform Q1.1b fallback) ---
\usepackage{fontspec}
\usepackage{xeCJK}
% Windows (local dev): SimSun/SimHei/FangSong
% Linux/macOS (CI): Noto Serif CJK SC / Noto Sans CJK SC (fonts-noto-cjk apt package)
% Last resort: AR PL UMing CN (fonts-arphic-uming) which is preinstalled in many texlive distros
\IfFontExistsTF{SimSun}{%
  \setCJKmainfont{SimSun}[BoldFont=SimHei]%
}{%
  \IfFontExistsTF{Noto Serif CJK SC}{%
    \setCJKmainfont{Noto Serif CJK SC}[BoldFont=Noto Serif CJK SC]%
  }{%
    \setCJKmainfont{AR PL UMing CN}[AutoFakeBold=2]%
  }%
}
\IfFontExistsTF{SimHei}{%
  \setCJKsansfont{SimHei}%
}{%
  \IfFontExistsTF{Noto Sans CJK SC}{%
    \setCJKsansfont{Noto Sans CJK SC}%
  }{%
    \setCJKsansfont{AR PL UKai CN}%
  }%
}
\IfFontExistsTF{FangSong}{%
  \setCJKmonofont{FangSong}%
}{%
  \IfFontExistsTF{Noto Sans Mono CJK SC}{%
    \setCJKmonofont{Noto Sans Mono CJK SC}%
  }{%
    \setCJKmonofont{AR PL UMing CN}%
  }%
}
\setmainfont{Times New Roman}[Ligatures=TeX]
\xeCJKsetup{CJKecglue={\hskip 0.15em plus 0.05em minus 0.05em}}

% --- Math ---
\usepackage{amsmath, amssymb, amsthm}

% --- Tables & Figures ---
\usepackage{booktabs}
\usepackage{graphicx}
\usepackage{longtable}
\usepackage{rotating}
\usepackage{caption}
\captionsetup{font=small,labelfont=bf,skip=4pt}

% --- Layout ---
\usepackage[margin=1in]{geometry}
\usepackage{setspace}
\onehalfspacing
\usepackage{hyperref}
\hypersetup{colorlinks=true,linkcolor=blue,citecolor=blue,urlcolor=blue}

% --- Bibliography ---
\usepackage[backend=biber,style=apa,natbib=true]{biblatex}

% --- Macros ---
\newcommand{\code}[1]{\texttt{#1}}
\usepackage[dvipsnames]{xcolor}
\newcommand{\NoteThis}[1]{\textbf{\textcolor{red}{[TODO: #1]}}}

% --- Theorem environments ---
\newtheorem{prop}{命题}
\newtheorem{hyp}{假设}


\addbibresource{refs.bib}

\title{媒体语调、金融科技战略与银行绩效：\\
       基于中国银行—年度新闻语料的证据}
\author{\NoteThis{作者姓名}}
\date{\today}

\begin{document}
\maketitle

\begin{abstract}
\NoteThis{最终摘要待定。}
本文基于约 37~GB 中国媒体金融科技报道构建银行—年度面板,
研究银行金融科技战略的媒体语调与银行经营结果之间的关系。
通过三家族大语言模型 (LLM) 集成进行情感标注,
我们同时度量金融科技报道的"主体性 (locus)"
(全文 vs.\ 实体上下文) 和"实质性 (substantiveness)"
(全文 vs.\ 实体上下文)。我们发现 \NoteThis{核心结论占位}\@。
\end{abstract}

\section{引言}\label{sec:intro}
\NoteThis{阐释问题: 媒体对银行金融科技战略的报道,
对制度理论 (Westphal \& Zajac 1994) 与数字化转型结果
(Gomber et al.\ 2018, \emph{JMIS}) 至关重要。借鉴 Hassan et al.\
(2019, \emph{QJE}) 公司层级政治风险方法作为我们的文本测度模板。}

\section{文献综述}\label{sec:lit}
\NoteThis{三条主线: (1)~媒体作为信息中介 (Tetlock 2007);
(2)~企业回应的象征 vs.\ 实质 (Westphal--Zajac 1994/1998 \emph{ASQ});
(3)~Bartik 工具变量设计 (Goldsmith-Pinkham et al.\ 2020, \emph{AER})。}

\section{方法}\label{sec:method}
\subsection{样本与数据}
银行—年度面板覆盖 \NoteThis{2017--2018 年共 N=164 个银行—年观测
(后续合并新一轮年报数据后扩展)}。原始文本输入为 XLSX 媒体报道,
经 \code{news\_process} 管线 (Phase~9b 基线) 处理。

\subsection{语调标注 (F1: 三家族 LLM 真三角验证)}
我们采用三个独立的 LLM 家族进行情感标注:
\code{gpt-5.4-nano} (OpenAI)\,、\code{gemini-3.1-flash-lite} (Google)
和 \code{deepseek-v4-flash} (DeepSeek)。
其中 OpenAI 与 Google 两家通过 yunwu.ai 聚合器以 OpenAI 兼容的
\code{/v1/chat/completions} API 调用; DeepSeek 直连
\code{api.deepseek.com}。所有调用设置 \code{temperature=0} 与
\code{seed=42}, 并使用字节级冻结的 prompt (版本 v8),
保证输出确定且可经缓存复现。
家族间一致性以 Cohen's $\kappa$ (计算前先把标签重映射为非负码)
与 Krippendorff's $\alpha$ (名义尺度、缺标感知、兼容二分类与三分类
标签集) 报告; 最终情感标签为家族多数投票, 平局标记为
\code{undefined} (而非 \code{neutral})。

除多数投票共识 (\code{llm\_vote\_tone\_full}, 即 spec 规定的
headline 处理变量) 外,
我们保留三家各自的语调序列 (\code{gpt54nano\_tone}、
\code{gemini31\_tone}、\code{deepseekv4\_tone}) 作为
\emph{互为稳健性、不分主次}的并列测度: 三个 per-model 序列并排报告,
各自进入恰好识别 (just-identified) 的独立设定, 绝不合并进同一个
过度识别系统 (同一批文章派生的工具共享共同冲击, 联合 Hansen $J$
检验无信息量)。由于各家序列的原始方差不同, 跨模型系数比较一律先做
$z$-score 标准化。per-model 设定族的族错误率 (family-wise error)
由 Romano--Wolf step-down 重抽样控制 (\code{rwolf}; 见
Table~A12), 而非以"三家都显著"的叙述性说法代替正式检验。

\subsection{篇章级编码 (F2: 中性标签去污染)}
若某文章 (a)~三家族 LLM 全部调用失败 (\code{llm\_call\_failure=1})
或 (b)~三票分裂为三方互不一致 (\code{llm\_undefined=1}),
则该文章情感编码为\emph{缺失}, 而非 \code{neutral}。
若银行—金融科技相关性在全文轨上的多数验证未通过
(\code{llm\_verified\_full=False}), 情感同样视为缺失。
此设定相比早期版本是有意调整: 旧版强制兜底为"中性",
导致 81.7\,\% 的中性标签实际是测量失败或验证拒绝。

\paragraph{全文轨独占的留篇门控。}
一篇文章是否被计为银行—金融科技报道, 仅由全文轨的验证结果
(\code{llm\_verified\_full}) 决定。实体上下文轨 (\code{ctx}) 与
银行—金科同句共现轨 (\code{ss}) 为\emph{以全文轨为条件}
(conditional-on-full) 的测度: 其行宇宙是全文轨留篇集的子集,
仅作为测量窗口稳健性轨, 不构成独立的留篇门控。相较早期
"任一轨验证通过即保留"的版本, 这一更严格的门控会机械性地压缩
文章计数及其派生份额 (\code{n\_articles}、\code{fintech\_share}、
\code{attention\_pct})。这些水平位移是"留篇门控与 headline
处理轨对齐"的定义性后果, 并非数据错误; 且方向偏保守
(剔除的是仅有更窄窗口才能验证的边缘文章)。

\subsection{零观测语义 (F3: NaN 主路径)}
对于无金融科技新闻覆盖的银行—年 (\code{n\_articles=0}),
所有基于余弦相似度的特征 (\code{auth\_sim\_*}、\code{benefit\_*}、
\code{article\_sim\_*})、比例变量
(\code{token\_adopt\_*}、\code{hidden\_adopt\_*}、
\code{subst\_adopt\_*})、连续强度分 (\code{fintech\_intensity\_mean})
均编码为未定义 (\code{NaN}), 而非零。
我们遵循 Allison (2001) 和 Heckman (1979) 把这类设计性缺失视为缺失值,
回归中采用列删除处理。计数与求和型变量在零观测时保留 $0$,
作为数学上正确的零观测值。

\subsection{Bartik IV 修剪 (F4)}
早期版本面板含约 30 个候选 Bartik 工具变量, 其中 28 个被 v6 删除,
原因是它们违反 Goldsmith-Pinkham et al.\ (2020, AER) 所要求的外生性。
保留的工具变量为 \code{bartik\_type\_iv\_full} 与
\code{prov\_leave1out\_tone\_full}。

\subsection{嵌入类语义测度 (双主轨、互为稳健性; F5)}
所有基于嵌入的特征 (金融科技强度、政策一致性、语义陈旧度及相关
相似度测度) 均在\emph{来自独立厂商的两个同为主测度的嵌入模型}下
并行计算、并排报告、互为稳健性: 阿里
\code{Qwen3-Embedding-8B} (4{,}096 维, 经硅基流动 API;
canonical 无后缀列) 与 OpenAI \code{text-embedding-3-large}
(3{,}072 维; \code{\_lg} 后缀列)。两轨遵循严格的列对偶
(任一嵌入特征要么两轨齐备、要么两轨皆无), 各轨进入各自恰好识别的
独立设定, 绝不合并进同一个过度识别系统。跨厂商独立性
(阿里 vs.\ OpenAI) 可防范单一厂商的表示偏差驱动任何结论。

\emph{换模披露。}早期管线版本的 canonical 槽使用 Google Gemini
嵌入模型; 自管线 v3.15.0 (2026 年 6 月) 起 canonical 槽换为阿里
\code{Qwen3-Embedding-8B}, 全部嵌入派生特征经受控全量重嵌再生,
且嵌入空间的语义效度由测试套件中的跨厂商 Spearman 秩相关闸把关
(两条 co-primary 嵌入空间两两相关 $\rho \geq 0.80$、$\geq 500$ 对)。
情感标注的 LLM 集成
(见方法节) 不受此变更影响: 嵌入模型与情感 LLM
是管线中互不相交的两类组件。

\section{数据描述}\label{sec:data}
\NoteThis{Table~1 描述统计参见
\code{news\_process/papers/bank\_fintech\_tone/tables.py:build\_table1}。}

\section{基准回归}\label{sec:main}
\NoteThis{Table 2 (M1--M6 渐进控制) 与 Figure 2 (系数图)
源自 \code{output/tables/table2\_main\_m1m6.csv}。}

\section{机制检验}\label{sec:mech}
\NoteThis{Table 3 (机制: 技术应用 vs.\ 实质--象征差异)
源自 \code{output/tables/table3\_mechanism.csv}。}

\section{异质性}\label{sec:het}
\NoteThis{Table 4 (银行规模 $\times$ 前沿性 $\times$ 地理) 源自
\code{output/tables/table4\_heterogeneity.csv}。}

\section{稳健性}\label{sec:robust}
\NoteThis{Table 5 (安慰剂、留一法、wild bootstrap) 加上 1016 行
falsification 模板
(\code{papers\_assets/bank\_fintech\_tone/do/falsification\_template.do.tmpl})。}

跨模型稳健性: headline 设定分别以三家 LLM 各自的语调测度重估
(BERT 侧同理, 以三个中文金融域 BERT 模型各自的有符号净语调
net\_tone 重估), per-model 系数一律 $z$-score 标准化,
族错误率由 Romano--Wolf step-down 重抽样控制 (Table~A12)。
模型间一致性诊断 (Cohen's $\kappa$, 非负码重映射后计算;
Krippendorff's $\alpha$, 缺标感知) 随表一并报告。

\section{结论}\label{sec:concl}
\NoteThis{最终讨论与政策含义。}

\printbibliography

\end{document}
